\section{Only the last-capacity edge can be flat before the first positive-derivative mode}
This statement concerns SINGLE normalized formal arms. Capacities after
the initial ones may occur in arbitrary order, provided they are at least
two. It is independent of the original one-descent conjecture.

Let $D_0=D_1=0$, $D_n=(F_n-F_{n-1})^+$ for $n\ge2$, and let
$\ell$ be its FIRST global mode. For a nonrectangular arm write $t$
for its final capacity and $M=D_\ell$.

\begin{theorem}[A single possible flat edge]
For every $3\le n\le\ell$,
\[
      D_n>D_{n-1}\quad\hbox{unless }n=t+1.
\]
There is at most one flat edge on this rising prefix. Consequently
\[
                             \boxed{\ell\le M+2.}       \tag{1}
\]
The height bound is sharp, for example on $1^\alpha2^3$ with
$\alpha\ge3$, where the first derivative mode is four and its height two.
\end{theorem}
\begin{proof}
Induct on the non-one count. A one-capacity arm has positive derivative
an interval of ones starting at rank two, whose FIRST mode is two;
the asserted rising-edge interval is empty.

For a further append $t\ge2$ let $d$ be the OLD positive derivative,
with first mode $m$, last positive rank $R$, and $d_2=1$.
By induction there is at most ONE flat edge through $m$; all other
edges there strictly increase. Put $W=S_td$.
For $3\le n\le m+1$ its increment is
\[
                         \Delta W_n=d_{n-1}-d_{n-t-1}>0. \tag{2}
\]
If the leaving index is below two, the entering coefficient is positive
and the leaving one zero. Otherwise the two indices are at least two
apart on the old rising prefix, which contains at most one flat edge;
at least one intervening edge is strictly positive. This proves (2).

For $n>m+1$ while the leaving index is at most $m$, the entering
coefficient is nonincreasing and the leaving coefficient nondecreasing.
Thus the increments in (2) are nonincreasing across that crossing
range. Once the leaving index is past $m$, each increment is nonpositive.
A zero crossing increment can therefore never be followed by a positive
one. Combining this with (2), $W$ strictly increases all the way to its
FIRST mode $m_W$; there is no earlier flat edge.

The exact positive-derivative splice is
\[
                          P=z^2[t-1]+W.
\]
The additional interval is nonincreasing from rank two, and $m_W\ge3$.
Hence $P$ is nonincreasing from $m_W$, so its FIRST mode $m_P\le m_W$.
For $3\le n\le m_P$, every increment of $P$ is the strictly positive
increment of $W$, except at the one edge $n=t+1$, where one unit is
subtracted. The previously established unimodality of $P$ guarantees
that this exceptional increment cannot be negative before its first mode.
Thus it is the only possible flat edge.

Finally $P$ is nonincreasing from $R+1$: its entering old derivative
is zero afterward and the additional interval has no positive increment
there. The actual NEW derivative equals $P$ through at least $R+1$,
and the known terminal splice is nonincreasing afterward. Its first
mode therefore agrees with that of $P$, and its entire rising prefix
has exactly the asserted strictness property. This closes induction.

For (1), $D_2=1$ and the $\ell-2$ integer increments entering ranks
three through $\ell$ are nonnegative, with at most one zero.
If $\ell\ge3$, they sum to at least $\ell-3$, so $M\ge\ell-2$;
the case $\ell=2$ is immediate. For the sharpness family the actual
derivative is $z^2+z^3+2z^4+2z^5+z^6$ for $\alpha\ge3$;
its first mode is four and height two.
\end{proof}

\paragraph{A smaller alternative finite region for the triangular-source proof.}
In a peak-crossing TC comparison the earlier unbounded payment gives
$\Delta\widehat E_n\ge M+n-\ell-3\ge n-5$ by (1).
At OLD $\kappa\le7$, the new triangular budget is at most 35.
Thus any failure has $n\le39$. The simpler peak payment
$\Delta\widehat E_n\ge M-1$ also forces $M\le35$ in a failure.
This gives a smaller replacement for the earlier endpoint-45/height-128
finite region, covering ALL old modes at small first weight. The original
independent endpoint-45 proof remains valid. The NEW endpoint-39/height-35
grammar has 102630 states and 1623669 legal appends; all 18321840 required
necessary-guard comparisons satisfy the exact source budget. Its independent
rolling-window DFS reconstruction has exactly the same counts and complete
state set, compared byte for byte with the fixed-width state certificate.
Thus it supplies an alternative peak-crossing proof in the full TC theorem,
without the late-old-peak case split. It is not a new original UI theorem.

The written strict-window induction above is the unbounded proof.
The separate diagnostic scans all 657394 nonrectangular coefficient
prefix states in the earlier complete bounded-slope grammar and 15000
wide actual arms; 1333 possible interior flat edges all occur at the
final-capacity edge, with no failure. These finite checks are not the
proof and do not identify arbitrary formal words with actual unsorted hosts.
